\documentclass{article}
\usepackage[margin=1in, a4paper]{geometry}
\usepackage{amsmath, amssymb, amsfonts}
\usepackage{graphicx}
\usepackage{xcolor}
\usepackage{listings}
\usepackage{booktabs}
\usepackage[utf8]{inputenc}
\usepackage[T1]{fontenc}
\usepackage{hyperref}
\hypersetup{colorlinks=true, urlcolor=blue, linkcolor=blue}
\usepackage{enumitem}
\usepackage{float}

\definecolor{codegreen}{rgb}{0,0.6,0}
\definecolor{codegray}{rgb}{0.5,0.5,0.5}
\definecolor{codepurple}{rgb}{0.58,0,0.82}
\definecolor{backcolour}{rgb}{0.99,0.99,0.99}
% Professional color palette for academic publishing
\definecolor{myblue}{cmyk}{0.8,0.4,0,0.1}
\definecolor{mygreen}{cmyk}{0.7,0,0.5,0.2}
\definecolor{myorange}{cmyk}{0,0.5,0.8,0.1}
\definecolor{mygray}{cmyk}{0,0,0,0.4}
\definecolor{arrowcolor}{cmyk}{0.9,0.5,0,0.3}
\definecolor{nodecolor}{cmyk}{0.6,0.2,0,0.05}
\definecolor{framecolor}{cmyk}{0.4,0.1,0,0.1}

\lstdefinestyle{mystyle}{
    backgroundcolor=\color{backcolour},   
    commentstyle=\color{codegreen},
    keywordstyle=\color{magenta},
    numberstyle=\tiny\color{codegray},
    stringstyle=\color{codepurple},
    basicstyle=\ttfamily\footnotesize,
    breakatwhitespace=false,         
    breaklines=true,                 
    captionpos=b,                    
    keepspaces=true,                 
    numbers=left,                    
    numbersep=5pt,                  
    showspaces=false,                
    showstringspaces=false,
    showtabs=false,                  
    tabsize=2
}
\lstset{style=mystyle}

\title{The Safety Stock \& Reorder Point (Q,r) Policy Problem}
\author{The Logistics \& Supply Chain Problem-Solving Playbook}
\date{}

\begin{document}
\maketitle

\section{The Safety Stock \& Reorder Point (Q,r) Policy Problem}

The (Q,r) inventory policy represents one of the most fundamental and widely implemented inventory management frameworks in supply chain operations. This policy defines when to order (\textit{r}, the reorder point) and how much to order (\textit{Q}, the order quantity) while maintaining appropriate safety stock levels to buffer against demand and lead time uncertainties. The challenge lies in optimally balancing inventory holding costs against stockout risks while considering the stochastic nature of demand and supply lead times.

\subsection{The Scenario: Managing Critical Components at Phoenix Distribution Center}

MedSupply Corp operates a major pharmaceutical distribution center in Phoenix, Arizona, serving hospitals and clinics across the southwestern United States. The facility manages over 15,000 different medical products, from basic supplies like syringes and bandages to critical life-saving medications with strict expiration dates and regulatory requirements.

Consider their management of a critical cardiac medication, CardioStabil, which has the following characteristics:
\begin{itemize}
\item Annual demand: 12,000 units with coefficient of variation of 0.3
\item Lead time: 14 days with standard deviation of 3 days
\item Unit cost: \$450 per vial
\item Holding cost rate: 25\% per year
\item Ordering cost: \$200 per order
\item Stockout cost: \$2,000 per unit (including expedited shipping and potential patient impact)
\item Required service level: 99.5\% (pharmaceutical industry standard)
\end{itemize}

The distribution center operates 365 days per year and faces the challenge of determining optimal safety stock levels and reorder points while minimizing total inventory costs. Demand exhibits both seasonal patterns and random variability, while supplier lead times can vary significantly due to manufacturing complexities and regulatory inspections.

\subsection{Tier 1: The Pen \& Paper Method (Mathematical Formulation)}

The mathematical foundation of the (Q,r) policy can be formulated as a network flow optimization problem where inventory flows through time and states, with safety stock serving as a buffer against uncertainty.

\textbf{Network Flow Formulation:}

Let us model the inventory system as a directed graph $G = (V, E)$ where:
\begin{itemize}
\item $V$ represents inventory states and time periods
\item $E$ represents possible transitions (ordering decisions and demand realizations)
\end{itemize}

\textbf{Sets and Parameters:}
\begin{align}
T &= \{1, 2, \ldots, T_{max}\} \quad \text{Time periods} \\
S &= \{0, 1, \ldots, S_{max}\} \quad \text{Inventory states} \\
\mu &= \text{Average demand per period} \\
\sigma_D &= \text{Standard deviation of demand} \\
L &= \text{Average lead time} \\
\sigma_L &= \text{Standard deviation of lead time} \\
h &= \text{Holding cost per unit per period} \\
K &= \text{Fixed ordering cost} \\
p &= \text{Stockout penalty cost per unit}
\end{align}

\textbf{Decision Variables:}
\begin{align}
x_{st} &= \text{Inventory level in state } s \text{ at time } t \\
y_{t} &= \text{Binary variable: 1 if order placed at time } t, 0 \text{ otherwise} \\
Q_{t} &= \text{Order quantity at time } t \\
r &= \text{Reorder point}
\end{align}

\textbf{Objective Function:}
Minimize the expected total cost over the planning horizon:
\begin{align}
\min \sum_{t \in T} \left[ h \cdot E[x_{st}] + K \cdot E[y_t] + p \cdot E[\max(0, D_t - x_{st})] \right]
\end{align}

\textbf{Network Flow Constraints:}
\begin{align}
x_{s,t+1} &= x_{st} + Q_t - D_t \quad \forall s \in S, t \in T \\
\sum_{s \in S} x_{st} &= 1 \quad \forall t \in T \quad \text{(Flow conservation)} \\
y_t &= 1 \text{ if } x_{st} \leq r \quad \forall t \in T \\
Q_t &\geq 0 \quad \forall t \in T \\
r, Q &\geq 0
\end{align}

\textbf{Safety Stock Calculation:}
The safety stock under demand and lead time uncertainty follows:
\begin{align}
SS &= z_{\alpha} \sqrt{L \sigma_D^2 + \mu^2 \sigma_L^2}
\end{align}

where $z_{\alpha}$ is the safety factor corresponding to the desired service level $\alpha$.

\textbf{Reorder Point Formulation:}
\begin{align}
r &= \mu \cdot L + SS = \mu \cdot L + z_{\alpha} \sqrt{L \sigma_D^2 + \mu^2 \sigma_L^2}
\end{align}

\begin{figure}[htbp]
\centering
\includegraphics[width=\textwidth]{../figures-png/line56.fig1.png}
\caption{Network Flow Model of (Q,r) Inventory Policy}
\end{figure}

\textbf{Concrete Example:}
Consider CardioStabil with the following parameters:
\begin{itemize}
\item Daily demand: $\mu = 33$ units, $\sigma_D = 10$ units
\item Lead time: $L = 14$ days, $\sigma_L = 3$ days
\item Service level: 99.5\% ($z_{0.995} = 2.576$)
\end{itemize}

\textbf{Step 1: Calculate Safety Stock}
\begin{align}
SS &= 2.576 \sqrt{14 \times 10^2 + 33^2 \times 3^2} \\
&= 2.576 \sqrt{1400 + 9801} \\
&= 2.576 \sqrt{11201} \\
&= 2.576 \times 105.84 \\
&= 272.6 \text{ units}
\end{align}

\textbf{Step 2: Calculate Reorder Point}
\begin{align}
r &= 33 \times 14 + 272.6 = 462 + 272.6 = 734.6 \approx 735 \text{ units}
\end{align}

\textbf{Step 3: Calculate Economic Order Quantity}
\begin{align}
Q^* &= \sqrt{\frac{2 \times 12000 \times 200}{450 \times 0.25}} = \sqrt{\frac{4,800,000}{112.5}} = \sqrt{42,667} = 206.6 \approx 207 \text{ units}
\end{align}

Therefore, the optimal policy is (Q = 207, r = 735) with safety stock of 273 units.

\subsection{Tier 2: The Classic Heuristic (Priority-Based Algorithm)}

The priority-based algorithm for (Q,r) policy optimization uses a systematic approach that prioritizes different cost components and service level requirements to iteratively converge on optimal parameter values.

\textbf{Algorithm Theory:}
The algorithm employs a priority queue structure where inventory decisions are prioritized based on criticality scores that combine cost ratios, service level requirements, and demand variability measures. This approach is particularly effective for multi-item inventory systems where computational efficiency is crucial.

\textbf{Time Complexity Analysis:}
The algorithm operates in $O(n \log n + k)$ time complexity, where $n$ is the number of inventory items and $k$ is the number of iterations for convergence. The logarithmic component comes from priority queue operations, while the linear component represents the iterative refinement process.

\lstinputlisting[language=Python, caption={Priority-Based (Q,r) Policy Algorithm}]{.././codes-python/line56.py1.py}

\begin{figure}[htbp]
\centering
\includegraphics[width=\textwidth]{../figures-png/line56.fig2.png}
\caption{Priority-Based Algorithm Flow for (Q,r) Policy Optimization}
\end{figure}

\textbf{Concrete Example Output:}
Running the CardioStabil example produces the following results:

\begin{verbatim}
Criticality Score: 178.33

Optimal Solution (Strategy: cost_focused):
Order Quantity (Q): 207 units
Reorder Point (r): 689 units  
Safety Stock: 227 units
Total Annual Cost: $67,440.25
\end{verbatim}

The algorithm identifies CardioStabil as a high-criticality item due to its high stockout costs relative to holding costs, resulting in a cost-focused strategy that balances inventory investment with service level requirements.

\subsection{Tier 3: The Advanced Algorithm (Particle Swarm Optimization)}

Particle Swarm Optimization (PSO) provides a powerful metaheuristic approach for optimizing (Q,r) policies, particularly effective when dealing with complex cost functions, multiple constraints, and non-linear relationships between policy parameters.

\textbf{PSO Theory for Inventory Optimization:}
In the context of (Q,r) policy optimization, each particle represents a potential solution with position vector $(Q_i, r_i)$ and velocity vector $(v_{Q_i}, v_{r_i})$. The swarm explores the solution space by balancing exploitation of known good solutions with exploration of new regions, making it particularly suitable for inventory problems with multiple local optima.

\textbf{Adaptive Parameter Strategy:}
This implementation uses adaptive parameters where inertia weight decreases linearly from 0.9 to 0.4, and acceleration coefficients adjust based on swarm convergence to prevent premature convergence while ensuring solution quality.

\lstinputlisting[language=Python, caption={Particle Swarm Optimization for (Q,r) Policy}]{.././codes-python/line56.py2.py}

\begin{figure}[htbp]
\centering
\includegraphics[width=\textwidth]{../figures-png/line56.fig3.png}
\caption{Particle Swarm Optimization Search Process}
\end{figure}

\textbf{Concrete Example Output:}
Running the PSO optimization for CardioStabil yields:

\begin{verbatim}
Converged at iteration 73

PSO Optimization Results for CardioStabil:
Optimal Q: 205 units
Optimal r: 712 units
Safety Stock: 250 units
Total Annual Cost: $66,890.45
Service Level Achieved: 0.9952
Converged in 74 iterations
\end{verbatim}

The PSO algorithm discovers a solution with slightly lower total cost than the heuristic approach while maintaining the required service level, demonstrating the effectiveness of metaheuristic optimization for complex inventory policy problems.

\subsection{Tier 4: The AI/ML/RL Augmentation Method (Deep Learning)}

Deep learning approaches transform (Q,r) policy optimization by learning complex patterns in demand, lead time variations, and cost relationships that traditional methods cannot capture. This neural network-based approach uses historical data to predict optimal policy parameters dynamically.

\textbf{Deep Learning Architecture:}
The system employs a multi-layer feedforward neural network with specialized layers for feature extraction, pattern recognition, and policy parameter prediction. The architecture includes dropout layers for regularization and batch normalization for stable training.

\lstinputlisting[language=Python, caption={Deep Learning for Dynamic (Q,r) Policy Optimization}]{.././codes-python/line56.py3.py}

\begin{figure}[htbp]
\centering
\includegraphics[width=\textwidth]{../figures-png/line56.fig4.png}
\caption{Deep Learning Architecture for Dynamic (Q,r) Policy}
\end{figure}

\textbf{Concrete Example Output:}
The deep learning model training and CardioStabil optimization produces:

\begin{verbatim}
Generating training data...
Training deep learning model...
Epoch 45/80: Early stopping triggered
Training completed with validation loss: 0.0023

Optimizing CardioStabil policy...

Deep Learning Results for CardioStabil:
Optimal Q: 203 units
Optimal r: 718 units
Safety Stock: 256 units
\end{verbatim}

The deep learning approach demonstrates the ability to learn complex patterns from diverse scenarios and provides policy recommendations that account for nuanced relationships between demand variability, cost structures, and market conditions that traditional analytical methods might miss.

\section{Practice Questions}

\textbf{Tier 1 Questions (Mathematical Formulation):}

1. A pharmaceutical distributor manages insulin pens with the following characteristics: annual demand of 8,400 units, daily demand standard deviation of 12 units, lead time of 10 days with standard deviation of 2 days, holding cost rate of 30\% annually, unit cost of \$85, ordering cost of \$150, and required service level of 99\%. Formulate the complete mathematical model and calculate the optimal safety stock and reorder point.

2. For a medical device with annual demand of 2,400 units, lead time of 21 days (CV = 0.25), daily demand CV of 0.4, and service level requirement of 99.5\%, derive the network flow constraints and calculate the minimum total inventory investment required.

\textbf{Tier 2 Questions (Heuristic Algorithms):}

3. Implement the priority-based algorithm for a portfolio of three medical products with the following data: Product A (annual demand: 15,000, high stockout cost), Product B (annual demand: 8,000, moderate variability), Product C (annual demand: 25,000, low service level requirement). Compare the criticality scores and resulting policies.

4. A surgical supply distributor faces seasonal demand with 40\% higher demand in winter months. Modify the heuristic algorithm to handle this seasonality and calculate the adjusted safety stock levels for a product with base annual demand of 6,000 units.

\textbf{Tier 3 Questions (Metaheuristic Optimization):}

5. Design a PSO experiment comparing swarm sizes of 20, 40, and 60 particles for optimizing a high-value cardiac medication with unit cost \$1,200, annual demand 3,600 units, and very high stockout penalties (\$5,000 per unit). Analyze convergence behavior and solution quality.

6. A specialty pharmacy manages 50 different medications with interdependent demand patterns (correlation coefficient 0.3). Extend the PSO algorithm to handle portfolio optimization and determine if centralized optimization provides cost benefits over individual product optimization.

\textbf{Tier 4 Questions (AI/ML Enhancement):}

7. The deep learning model shows prediction accuracy of 94\% on validation data. For a new product with limited historical data (only 6 months), design a transfer learning approach that leverages knowledge from similar products to predict optimal (Q,r) parameters.

8. A hospital network operates in three different market conditions: urban high-competition (fast-moving inventory), suburban moderate-competition (stable demand), and rural low-competition (slow-moving, high service requirements). Train separate neural network models for each market type and compare their policy recommendations for the same product characteristics.

\end{document}
